\clearpage

%==============================================================================
% Appendix Section: Participant Perceptions of Study Intent
%==============================================================================

\section{Participant Perceptions of Study Intent}
\label{sec:appendix_study_intent}

%------------------------------------------------------------------------------
% Table A.5: Robustness excluding respondents who understood the study's intent
%------------------------------------------------------------------------------
\begin{table}[htbp]
\centering
\caption{Treatment effects excluding respondents who understood the study's intent}
\label{tab:treatment_effects_demand}
\begin{tabular}{l ccccc}
\toprule
 & (1) & (2) & (3) & (4) & (5) \\
 & SE (trial) & Delta & Vacc Intent & Link Click & Vaccinated \\ \midrule
Industry            &      -7.980&      -3.412&      -0.001&      -0.004&      -0.012\\
                    &     (0.901)&     (1.072)&     (0.013)&     (0.006)&     (0.012)\\
Academic            &      -3.911&      -2.908&       0.031&      -0.002&       0.015\\
                    &     (0.928)&     (1.072)&     (0.013)&     (0.006)&     (0.013)\\
Personal            &      -4.563&       0.530&       0.018&      -0.004&       0.012\\
                    &     (0.938)&     (1.089)&     (0.013)&     (0.006)&     (0.013)\\
Control mean        &      20.672&      15.480&       0.072&       0.015&       0.055\\
N                   &       3,186&       3,186&       3,177&       3,186&       2,697\\
 \\
\bottomrule
\end{tabular}
\tabnote{\linewidth}{Notes: Each column is a separate OLS regression of the
indicated outcome on treatment arm indicators for Industry, Academic, and
Personal arms, with Control as the omitted category. Pre-specified controls
are included but not shown. Robust standard errors in parentheses. The sample
is restricted to main-sample respondents whose response to ``What do you think
this study was about?'' was not classified as indicating they understood the
study's intent (Figure~\ref{fig:debrief_what_about}; classification described
in code/classification\_prompt.md). See Table~\ref{tab:treat_effects} for
outcome definitions.}
\end{table}

\clearpage

%------------------------------------------------------------------------------
% Figure A.3: Classification of "what was this study about" responses
%------------------------------------------------------------------------------
\begin{figure}[htbp]
\centering
\includegraphics[width=0.8\textwidth]{../output/figures/debrief_what_about.png}
\caption{Classification of Open-Ended ``What Was This Study About?'' Responses}
\label{fig:debrief_what_about}
\fignote{\linewidth}{Notes: Bars show, for the main sample, the percent of
responses to ``What do you think this study was about?'' classified by
claude-haiku-4.5 (prompt in code/classification\_prompt.md) as mentioning each
category. Categories are not mutually exclusive. ``Understood study's intent''
indicates the response was classified as mentioning flu/vaccine, information,
and change, and at least one of attitudes, beliefs, or intentions.}
\end{figure}

%------------------------------------------------------------------------------
% Table A.6: Example responses by classification pattern
%------------------------------------------------------------------------------
\begin{table}[htbp]
\centering
\caption{Example Responses by Classification Pattern}
\label{tab:debrief_examples}
{\setlength{\extrarowheight}{3pt}
\begin{tabularx}{\textwidth}{>{\raggedright\arraybackslash}p{0.4\textwidth} >{\raggedright\arraybackslash}X}
\toprule
Classification pattern & Example response \\
\midrule
Mentions flu or vaccine only & ``Flu vaccines'' \\
Mentions feelings, attitudes, or opinions only & ``How people react to placebos'' \\
Mentions information or research only & ``Paying attention to the study by reading the information given by the researcher.'' \\
Mentions side effects or effectiveness only & ``I am not sure. Maybe taking risks.'' \\
Mentions vaccination intentions only & ``To see what decisions I will make.'' \\
Mentions changing behavior, attitudes, or beliefs only & ``To see if you could be influenced'' \\
Mentions flu/vaccine and feelings/attitudes & ``The goal of the study was to evaluate participants perceptions of flu-shots.'' \\
Mentions flu/vaccine and information & ``Information about vaccinations'' \\
Mentions flu/vaccine and side effects/effectiveness & ``To see if people think flu vaccines cause adverse effects'' \\
Mentions flu/vaccine, feelings/attitudes, and information & ``perception of information and shots'' \\
Mentions flu/vaccine and changing behavior, without information & ``encourage flu vaccine'' \\
Understood study's intent (flu/vaccine, information, change, and attitudes, side effects/effectiveness, or intentions) & ``To see if people would change their mind on getting the vaccine upon learning about a study that shows people getting the vaccine aren't more likely to experience symptoms than those who got the placebo.''
\\
\bottomrule
\end{tabularx}}
\tabnote{\linewidth}{Notes: Each row shows one participant's verbatim response
to ``What do you think this study was about?'', illustrating a pattern of
categories (defined in code/classification\_prompt.md) assigned by the
claude-haiku-4.5 classifier.}
\end{table}

\clearpage

%------------------------------------------------------------------------------
% Classification prompt (verbatim)
%------------------------------------------------------------------------------
\subsection*{Classification Prompt}
\label{sec:classification_prompt}

\begin{sloppypar}
\noindent The following prompt was used (via
\path{code/classify_open_responses.py}) to classify each open-ended
``What do you think this study was about?'' response with claude-haiku-4.5
(\path{code/classification_prompt.md}):
\end{sloppypar}

\vspace{0.5em}
\lstinputlisting[basicstyle=\ttfamily\small, breaklines=true, breakindent=0pt]{../code/classification_prompt.md}
